\chapter{Preliminary Silicon Work: Executions, Finite-Setting Evidence, and an Exploratory Wannier Decomposition}
\index{silicon!preliminary evidence|textbf}

\label{app:preliminary-silicon-work}

\chapterstatus{bounded synthesis of preliminary silicon work. The tutorial and
finite-setting records cited here are retained repository evidence with their
own statuses. The same-frame Wannier kinetic decomposition is a separately
identified external, read-only prototype whose analysis package is not part of
this manuscript branch. Nothing in this appendix establishes an accepted
production silicon parent, converged Wannier representation, effective mass,
scientific validation, uncertainty quantification, or authorization for new
calculator execution.}

\section{Purpose and evidence boundary}

The silicon program has progressed beyond an entirely prospective plan, but it
has not crossed the production-material acceptance boundary stated in
Chapter~\ref{ch:results-and-limitations}. This appendix gathers the preliminary
work in one place so that successful execution observations, finite-setting
numerical evidence, and exploratory represented-operator analysis are not
mistaken for one homogeneous result.

The evidence is divided into four layers:

\begin{enumerate}
  \item a legacy Quantum ESPRESSO 7.2 silicon tutorial reproduction and its
        associated band calculation;
  \item a bounded Quantum ESPRESSO 7.5 rerun of the same tutorial-scale SCF
        input and a mechanical cross-version comparison;
  \item a direct PBE/PseudoDojo cutoff--mesh bootstrap matrix with a read-only
        audit, provisional finite-setting selection, and an unresolved runtime
        warning boundary; and
  \item a read-only analysis of a separate QE~7.5--Wannier90~3.1.0 silicon run,
        used to prototype a same-frame decomposition
        $H^W=T^W+R^W$, interpolation, local multiband expansion, and
        finite-mesh block truncation.
\end{enumerate}

These layers do not share one physical parent. In particular, the tutorial
uses a legacy LDA pseudopotential, the cutoff--mesh matrix uses the frozen
PseudoDojo PBE norm-conserving artifact, and the exploratory Wannier source run
uses the legacy ultrasoft \path{Si.pbe-n-van.UPF}. Their energies and operators
must not be combined as though they were refinements of one calculation.

\begin{table}[htbp]
  \centering
  \small
  \caption{Status of the preliminary silicon evidence summarized here.}
  \label{tab:preliminary-si-status}
  \begin{tabular}{p{0.22\linewidth}p{0.31\linewidth}p{0.37\linewidth}}
    \toprule
    Layer & Supported statement & Unsupported promotion \\
    \midrule
    QE 7.2 tutorial & The identified SCF reproduced the bundled printed energy
    and iteration count; a separate band run completed and retained 28 points
    with eight eigenvalues each. & Production convergence, modern band-path
    adequacy, a band gap, a valley, an effective mass, or pseudopotential
    suitability. \\
    QE 7.5 tutorial comparison & The bounded rerun completed and the listed
    scalar and QEXSD observations were mechanically compared with QE 7.2. &
    General cross-version equivalence or validation of either physical parent. \\
    PBE/PseudoDojo finite-setting matrix & The retained nine-case matrix passed
    its integrity audit; declared adjacent-setting criteria support a
    provisional C48/K8 selection. & Infinite-basis convergence, cutoff--mesh
    interaction control, accepted lattice geometry, effective-mass convergence,
    or downstream production eligibility. \\
    Exploratory Wannier decomposition & The reported finite rank-four matrices
    obey the same-frame decomposition and bounded numerical checks for the
    identified external source bytes. & An accepted production Wannier operator,
    short-range tight-binding model, converged curvature, or silicon validation. \\
    \bottomrule
  \end{tabular}
\end{table}

\section{Source records}

Table~\ref{tab:preliminary-si-sources} identifies the records from which this
appendix is projected. Repository paths are relative to the repository root.
The complete package in the final row remains external. A bounded publication
extraction containing the primary result, selected finite-mesh records and figures,
and a copy of the external output manifest is retained locally; the omitted package
payloads remain external.

\begin{table}[htbp]
  \centering
  \small
  \caption{Primary records for the preliminary silicon synthesis.}
  \label{tab:preliminary-si-sources}
  \begin{tabular}{p{0.23\linewidth}p{0.67\linewidth}}
    \toprule
    Evidence layer & Primary record or locator \\
    \midrule
    QE 7.2 tutorial SCF &
    \path{calculations/bulk-silicon/qe-example01-si-scf-davidson/result.md} \\
    QE 7.2 tutorial bands &
    \path{calculations/bulk-silicon/qe-example01-si-bands-davidson/result.md} \\
    QE 7.5 smoke comparison &
    \path{calculations/bulk-silicon/qe-7.5-si-scf-smoke-comparison/execution-comparison.json} \\
    Finite-setting audit &
    \path{calculations/bulk-silicon/production-convergence-preflight/direct-results-audit.md} \\
    Provisional disposition &
    \path{calculations/bulk-silicon/production-convergence-preflight/direct-results-setting-disposition.json} \\
    Warning classification &
    \path{calculations/bulk-silicon/production-convergence-preflight/ieee-warning-class-separated-nonmpi-c48-classification.md} \\
    External Wannier prototype &
    Local bounded extraction:
    \path{calculations/ICMSEP2026/conference/paper_1/}\
    \path{silicon-wannier-kinetic-preview/}; complete external package identifier
    \path{silicon-example11-hminus-t-readonly-analysis-20261004T193610Z};
    \path{result.json} SHA-256
    \path{846564b1d716044c9ed0ac4ddc94a1d6b7e81019ae2235979fcadab5033cfcf0};
    copied output-manifest SHA-256
    \path{5f509ab1cd4fa998f4c4002eba364982e94ee239237dd1adc6e1e29f1094dfb5} \\
    \bottomrule
  \end{tabular}
\end{table}

\section{Legacy tutorial baseline}
\index{silicon!Quantum ESPRESSO tutorial}

The retained QE~7.2 Davidson SCF tutorial run exited with status zero, emitted
\path{JOB DONE.}, and reported convergence after six iterations. Its printed
total energy was
\begin{equation}
  E_{\mathrm{tutorial}}=-15.84452726\ \mathrm{Ry},
\end{equation}
matching the bundled example output at printed precision. This is an accepted
tutorial reproduction, not a production material result. Its primary record is
\path{calculations/bulk-silicon/qe-example01-si-scf-davidson/result.md}.

The associated band calculation also exited with status zero and emitted the
end-of-band-structure and completion markers. It retained 28 ordered
reciprocal-space points with eight printed eigenvalues at each point. Of the
224 printed values, 166 were identical to the bundled legacy output at printed
precision, and the largest observed printed difference was approximately
$10^{-4}\ \mathrm{eV}$. No comparison tolerance was accepted, so that count is
an observation rather than a numerical pass criterion. The coordinates form
the legacy tutorial path; they are not an accepted modern silicon path or an
effective-mass stencil.

The band result file retains older ``awaiting human review'' prose, while the
later lifecycle summary records the tutorial band task as closed
human-accepted. This appendix does not use that administrative difference to
strengthen the scientific claim: under either description, the result remains
bounded tutorial evidence only.

Both tutorial invocations emitted
\begin{center}
\small\path{IEEE_INVALID_FLAG IEEE_DIVIDE_BY_ZERO IEEE_OVERFLOW_FLAG IEEE_UNDERFLOW_FLAG}.
\end{center}
Successful termination does not classify those flags as harmless.

\section{QE 7.5 smoke comparison}

A separately authorized, isolated, single-process QE~7.5 invocation reused the
same tutorial-scale SCF input and legacy pseudopotential identity. It also
completed in six iterations with printed total energy
$-15.84452726\ \mathrm{Ry}$. The parsed QEXSD energies differed by
$-2.00030214614344\times10^{-10}\ \mathrm{Ha}$, or
$-4.00060429228688\times10^{-10}\ \mathrm{Ry}$, between the QE~7.5 and QE~7.2
records. The largest listed absolute QEXSD eigenvalue difference was
$2.3826225714618232\times10^{-11}\ \mathrm{Ha}$.

These values are a bounded compatibility observation for one input. Native
charge-density and wavefunction files were not byte-identical, and byte
identity was not required. The exercise also motivated fail-closed parser
support for the exact observed QEXSD format versions \path{23.03.10} and
\path{25.05.21}; it did not establish exhaustive compatibility with either
upstream schema. The exact comparison and its limits are retained in
\path{calculations/bulk-silicon/qe-7.5-si-scf-smoke-comparison/execution-comparison.json}.
The same four-class IEEE report recurred.

\section{PBE/PseudoDojo finite-setting matrix}
\index{silicon!finite-setting evidence}

The physically relevant bootstrap branch is distinct from the tutorial. It
uses the frozen scalar-relativistic, non-SOC, non-spin-polarized PBE parent and
the exact PseudoDojo norm-conserving Si artifact required by the physical and
numerical specifications. The direct campaign contains six cutoff labels
\path{C30}, \path{C36}, \path{C42}, \path{C48}, \path{C54}, and \path{C60},
and four shifted-mesh labels \path{K6}, \path{K8}, \path{K10}, and
\path{K12}. The logical \path{K8} case reuses \path{C48}. Nine unique cases
therefore supplied nine SCFs and nine linked diagnostic NSCFs.

The later read-only audit recomputed the compact identities, checked all 18
input copies, matched 41 retained text-output and receipt identities,
independently parsed all nine cases, and reconstructed the eight adjacent
comparisons. Every invocation had one \path{JOB DONE.} marker and no retry. The
audit explicitly did not reread or duplicate the dense restart and wavefunction
payloads merely to restate compact observables.

The frozen adjacent-setting criteria were $10^{-5}\ \mathrm{Ry}$ per atom for
total energy, $0.05\ \mathrm{kbar}$ for pressure and each stress component,
and $1\ \mathrm{meV}$ for the retained fixed-point band and gap probes. The
first two cutoff comparisons failed at least the energy, pressure, and stress
criteria. Comparisons C42$\rightarrow$C48, C48$\rightarrow$C54, and
C54$\rightarrow$C60 met all listed criteria. All three mesh comparisons,
K6$\rightarrow$K8, K8$\rightarrow$K10, and K10$\rightarrow$K12, also met the
listed criteria at the printed resolution.

A later human disposition selected the smallest setting with passing lower and
upper finite-setting guards:
\begin{equation}
  E_{\mathrm{cut}}^{\psi}=48\ \mathrm{Ry},\qquad
  E_{\mathrm{cut}}^{\rho}=192\ \mathrm{Ry},\qquad
  \mathcal K_{\mathrm{SCF}}=8\times8\times8\ \text{shifted}.
\end{equation}
The selection is provisional. Sequential one-variable scans do not bound the
cutoff--mesh interaction, and the diagnostic points do not locate the
conduction valley or test curvature. No infinite-basis estimate follows.

\subsection{What the finite-setting data do and do not show}

The extracted per-atom energy changed over the cutoff sequence:
\begin{equation}
  E_{\mathrm{C30}}=-8.461848470\ \mathrm{Ry},\qquad
  E_{\mathrm{C60}}=-8.461975385\ \mathrm{Ry}.
\end{equation}
The fixed-point probes were numerically much less sensitive over the upper
cutoff sequence, but their
limited coordinates and four-decimal electron-volt printing make them
unsuitable substitutes for the later valley and effective-mass calculations.
Pressure remained near $37.8\ \mathrm{kbar}$ because the provisional geometry
was not the accepted zero-pressure PBE-relaxed lattice required by the physical
specification. Thus the matrix supports finite-setting behavior at a fixed
provisional geometry; it does not accept that geometry.

The authoritative synthesis is
\path{calculations/bulk-silicon/production-convergence-preflight/direct-results-audit.md}.
The selection and warning state are recorded separately in
\path{direct-results-setting-disposition.json}; a digest or successful process
exit is not a substitute for those semantic dispositions.

\section{Runtime warning diagnosis and remaining block}

Every SCF and diagnostic NSCF in the finite-setting matrix reported the same
four IEEE classes. Subsequent separately authorized diagnostics localized more
than the original text report alone could show:

\begin{enumerate}
  \item one first instrumented trap occurred in the external
        MPI/PMIx/hwloc/OpenCL startup path;
  \item after OpenCL discovery was disabled, a later floating-point division
        trapped below \path{ZHEGVX} during initial-wavefunction subspace
        diagonalization;
  \item that later operation appeared with both Apple Accelerate and OpenBLAS;
  \item a true non-MPI OpenBLAS build reproduced the matching OpenBLAS
        \path{ztrsv_CUN+388} operation; and
  \item a divide-by-zero-only non-MPI build reproduced the exact instruction,
        classifying that matching first later operation as an IEEE
        divide-by-zero exception.
\end{enumerate}

This localization does not recover the operands, explain the mathematical
cause, establish whether the produced value affects retained observables,
cover underflow, or identify every later invalid or overflow flag. The
provisional setting therefore remains blocked from downstream protected
execution pending an observable-effect and remaining-warning disposition. No
new calculator execution is authorized by this appendix.

\section{Exploratory same-frame Wannier kinetic decomposition}
\index{silicon!Wannier kinetic decomposition}

A separate local analysis consumed an externally retained QE~7.5--Wannier90
3.1.0 source run without executing either calculator. The source-run identifier
is
\begin{center}
\small\path{silicon-example11-hminus-t-20261004T171505Z}.
\end{center}
The analysis package identifies itself as
\path{silicon-example11-hminus-t-readonly-analysis-20261004T193610Z}. Its
\path{result.json} SHA-256 is
\begin{center}
\small\path{846564b1d716044c9ed0ac4ddc94a1d6b7e81019ae2235979fcadab5033cfcf0}.
\end{center}
The complete package remains in a separate development worktree and is not part
of the present manuscript branch. A bounded publication subset---the primary
result, truncation record, truncation figure, external output manifest, tutorial
source record, executed Wannier90 input and diff, and local checksums---is
retained under
\path{calculations/ICMSEP2026/conference/paper_1/}\
\path{silicon-wannier-kinetic-preview/}. Its hashes bind bytes observed by the
analysis, but the analysis did not consume a pre-existing authenticated
production-lineage manifest. Its runtime used content-bound working-tree source
over a base commit that alone is insufficient to reconstruct the
implementation. The numerical values below are therefore tutorial-scale
prototype observations, not promoted production evidence.

The numerical design is the official Wannier90~3.1.0
\path{examples/example11} silicon tutorial, not an arbitrary choice of
exploratory settings. A post-hoc byte-identity check confirms that the source
run's four original files match the retained upstream tutorial records. The
executed \path{silicon.win} is a derived copy adding only
\path{write_hr=true} and \path{write_u_matrices=true} output-artifact requests;
it changes no tutorial numerical setting. The tutorial uses the ultrasoft PBE
pseudopotential reference
\path{Si.pbe-n-van.UPF}, a 25-Ry wavefunction cutoff, an unshifted
$10\times10\times10$ SCF mesh, a $4\times4\times4$ Wannier mesh, 12 parent
bands, four retained Wannier functions, and $6.5\ \mathrm{eV}$ frozen- and
outer-window ceilings. This makes it an appropriate interface and analysis
preview, but not a production design: it cannot be joined to the PseudoDojo
C48/K8 branch by name, material, or software version alone. The tutorial's
familiar smooth band diagrams test interpolation of the total represented
Hamiltonian; their visual quality does not certify separate locality or
short-range compressibility of the represented kinetic and non-kinetic
components.

\subsection{Same-frame construction}

At each of the 64 retained reciprocal samples, the prototype composes the
rectangular disentanglement matrix and square Wannier gauge as
\begin{equation}
  W(\mathbf k)=U_{\mathrm{dis}}(\mathbf k)U(\mathbf k).
\end{equation}
Let $C(\mathbf k)$ contain the recovered pseudo-wavefunction plane-wave
coefficients and let $\boldsymbol\epsilon(\mathbf k)$ contain the parent band
energies. The represented operators are
\begin{align}
  H^W(\mathbf k)
    &=W(\mathbf k)^\dagger
      \operatorname{diag}\!\bigl(\boldsymbol\epsilon(\mathbf k)\bigr)
      W(\mathbf k),\\
  T^W(\mathbf k)
    &=W(\mathbf k)^\dagger C(\mathbf k)^\dagger
      \operatorname{diag}\!\left(\frac{|\mathbf k+\mathbf G|^2}{2}\right)
      C(\mathbf k)W(\mathbf k),\\
  R^W(\mathbf k)&=H^W(\mathbf k)-T^W(\mathbf k).
\end{align}
The kinetic diagonal is evaluated in Hartree before conversion to
electron-volts. The remainder $R^W$ is a finite represented non-kinetic
remainder; it is not identified with a continuous local scalar potential. It
may contain represented local, nonlocal-pseudopotential, Hartree,
exchange--correlation, and other parent-Hamiltonian contributions.

The source-backed convention audit reports a maximum UNK Fourier round-trip
defect of $2.16\times10^{-15}$ and a maximum Parseval defect of
$8.88\times10^{-16}$. The 25-Ry geometric support contains 537--568 plane
waves over the mesh. The even $24\times24\times24$ FFT grid has signed-Nyquist
ambiguity, so Nyquist-plane coordinates were excluded instead of assigned by
preference; their largest observed band weight was
$1.26\times10^{-32}$ for this source run.

Because the pseudopotential is ultrasoft, the parent eigenproblem has a
generalized overlap. The ordinary pseudo-coefficient norms, observed between
$0.98598$ and $1.02498$, were not forced to one. The UNK payload does not
contain the beta-projector augmentation required to reconstruct the overlap
matrix independently. The generalized-overlap interpretation is therefore
source-backed but not numerically reconstructed from UNK bytes alone.

\subsection{Reported finite-representation diagnostics}

The prototype reports reciprocal decomposition and Hermiticity defects at
approximately $10^{-15}\ \mathrm{eV}$ and Fourier round-trip defects below
$7\times10^{-15}\ \mathrm{eV}$. The independent comparison with native
\path{hr.dat} differs by $1.2763350779\times10^{-5}\ \mathrm{eV}$ in maximum
Frobenius norm on the retained construction and by approximately
$1.06\times10^{-5}\ \mathrm{eV}$ along the interpolated source-input path.
That cross-route mismatch is not folded into the same-route algebraic closure.

For each Hamiltonian eigenstate $v_n(\mathbf k)$, the reported expectation
identity is
\begin{equation}
  E_n(\mathbf k)
  =v_n^\dagger T^Wv_n+v_n^\dagger R^Wv_n,
\end{equation}
with maximum closure defect $5.33\times10^{-15}\ \mathrm{eV}$ over 256
state samples. This does not define separate kinetic and remainder band
structures: $T^W$ and $R^W$ are generally off diagonal in the Hamiltonian
eigenbasis, and their individual expectations inside an exactly degenerate
subspace depend on the chosen eigenvectors.

\subsection{Local rank-three expansion at \texorpdfstring{$\Gamma$}{Gamma}}

The reported represented Hamiltonian at $\Gamma$ contains one lower state and
an upper numerical triplet,
\begin{equation}
  \operatorname{eig} H^W(\Gamma)
  =(-5.6982916462,\ 6.3868784657,\ 6.3868784667,\ 6.3868784673)\ \mathrm{eV}.
\end{equation}
The triplet splitting is $1.57\times10^{-9}\ \mathrm{eV}$, below the declared
$10^{-8}\ \mathrm{eV}$ degeneracy tolerance. The analysis therefore uses a
rank-three selected-space quadratic tensor rather than differentiating three
independently ordered scalar bands through the degeneracy.

Analytic Cartesian derivatives of the Wigner--Seitz interpolation were checked
by centered finite differences at seven step sizes. Reported gradient and
Hessian convergence orders were approximately two for $H^W$, $T^W$, and
$R^W$. The quadratic model's largest ordered-eigenvalue error over the tested
directions and radii increased from $2.67\times10^{-10}\ \mathrm{eV}$ at the
smallest reported offset to $4.29\times10^{-5}\ \mathrm{eV}$ at the largest.
This checks the local expansion of the retained interpolation; it does not
check reciprocal-mesh or parent-model convergence.

The remote-state quadratic term contains kinetic--kinetic,
remainder--remainder, and kinetic--remainder products. The reported cross-term
norm reaches $8.864\ \mathrm{eV\,bohr^2}$. Separate additive ``kinetic'' and
``remainder'' effective masses are therefore not defined by this construction.
The prototype deliberately reports directional multiband curvatures without
converting them to silicon effective masses.

\subsection{Finite-mesh block truncation}

The 64 reciprocal matrices define exactly 64 centered lattice representatives.
The prototype tested inversion-closed radius truncations and reconstructed only
the same retained reciprocal mesh. Selected maximum relative matrix errors are
summarized in Table~\ref{tab:preliminary-si-truncation}.

\begin{table}[htbp]
  \centering
  \small
  \caption{Reported tutorial-scale finite-mesh truncation errors. The final column
  gives the maximum ordered-eigenvalue error for $H^W/T^W/R^W$.}
  \label{tab:preliminary-si-truncation}
  \begin{tabular}{rrrrrr}
    \toprule
    Blocks & Radius (bohr) & $H^W$ & $T^W$ & $R^W$ & Error (eV) \\
    \midrule
     1 &  0.000 & 0.594   & 0.376   & 0.206   & $4.04/6.35/2.51$ \\
    13 &  7.212 & 0.0665  & 0.0688  & 0.0641  & $0.558/1.467/0.942$ \\
    57 & 24.985 & 0.00196 & 0.00332 & 0.00421 & $0.0207/0.0843/0.0636$ \\
    63 & 29.738 & 0.000594& 0.00112 & 0.00127 & $0.00426/0.0202/0.0159$ \\
    64 & 35.334 & roundoff& roundoff& roundoff& roundoff \\
    \bottomrule
  \end{tabular}
\end{table}

The data do not support an aggressive short-range truncation. Even 57 retained
blocks permit a maximum kinetic eigenvalue error of about $84\ \mathrm{meV}$.
Thus a convincing total-band diagram and accurate complete finite
reconstruction can coexist with componentwise truncation errors. This is a
negative result about compression on one small periodic mesh, not an
infinite-lattice localization conclusion and not proof that a separately
fitted short-range model is impossible.

\section{What has not yet been established}

The preliminary work does not supply any of the following:

\begin{itemize}
  \item an accepted zero-pressure PBE-relaxed lattice and production SCF parent;
  \item a resolved observable-effect disposition for the recurring IEEE flags;
  \item a cutoff--mesh interaction check or an infinite-setting error bound;
  \item authenticated common lineage connecting the provisional PseudoDojo
        parent to the exploratory ultrasoft Wannier source run;
  \item an accepted Wannier rank, projection family, windows, gauge policy,
        mesh-convergence result, or independent withheld interpolation test for
        the production branch;
  \item a converged conduction-valley coordinate, longitudinal or transverse
        electron effective mass, or compatible SOC/spinor valence analysis;
  \item an aligned $sp^3s^{\ast}$ operator fit, model-class selection, or route
        comparison; or
  \item scientific validation or propagated uncertainty for a declared use.
\end{itemize}

The next scientifically meaningful sequence remains the one defined by the
production program: close the warning and finite-setting interaction boundary,
obtain and accept the zero-pressure lattice, construct the production SCF
parent, run purpose-specific path and local-valley children, and only then
freeze a production Wannier design. The exploratory decomposition is useful as
a prototype of representation and diagnostic machinery, but it cannot bypass
those gates.
